\documentclass[12pt,letterpaper]{article}
\usepackage{amsmath,amssymb,amsthm}
\newtheorem{theorem}{Theorem}
\newtheorem{proposition}[theorem]{Proposition}
\begin{document}
\section*{Supplemental mathematical proofs}

\section*{Appendix overview}

This appendix contains the full mathematical details that are referenced in the main text. The appendix uses sections only, as requested. Each proof begins with a one sentence interpretation so that the formal derivation stays connected to the main insight.

\section*{Proof of Proposition~\ref{prop:cluster}}\label{app:proofcluster}

Interpretation. This proof shows that the border score is a direct measure of expected local burst length.

\begin{proof}
Under the renewal cluster model, a cluster continues after each active occurrence with probability $\rho_w$ and stops with probability $1-\rho_w$. Hence
\begin{equation}
\mathbb{P}(C_w = k) = (1-\rho_w)\rho_w^{k-1},
\qquad k = 1,2,\ldots
\label{eq:appgeom}
\end{equation}
which is the geometric distribution on $\{1,2,\ldots\}$. Therefore
\begin{equation}
\mathbb{E}[C_w]
=
\sum_{k=1}^{\infty} k (1-\rho_w)\rho_w^{k-1}
=
\frac{1}{1-\rho_w}.
\label{eq:appgeommean}
\end{equation}
Using the definition $\beta_w=\rho_w/(1-\rho_w)$ yields
\begin{equation}
\frac{1}{1-\rho_w} = 1 + \frac{\rho_w}{1-\rho_w} = 1+\beta_w.
\label{eq:appbetaidentity}
\end{equation}
This completes the proof.
\end{proof}

\section*{Proof of Proposition~\ref{prop:variance}}\label{app:proofvariance}

Interpretation. This proof shows that valid self overlap creates positive local dependence and therefore larger variance of token counts.

\begin{proof}
For brevity write $I_t=I_t(w;X)$. Since $N_w=\sum_{t=1}^{T-m+1} I_t$, the variance decomposition gives
\begin{equation}
\mathrm{Var}(N_w)
=
\sum_{t=1}^{T-m+1}\mathrm{Var}(I_t)
+
2\sum_{1\leq s<t\leq T-m+1}\mathrm{Cov}(I_s,I_t),
\label{eq:appvardecomp}
\end{equation}
which is Equation~\eqref{eq:varcountbasic}. Stationarity implies $\mathbb{P}(I_t=1)=p_w$ and $\mathrm{Var}(I_t)=p_w(1-p_w)$.

Now fix a shift $\ell$ with $1\leq \ell < m$. If $\ell \notin S(w)$, then two occurrences of $w$ at positions $t$ and $t+\ell$ are incompatible because the required prefix and suffix letters disagree. Hence
\begin{equation}
\mathbb{P}(I_t=1,I_{t+\ell}=1)=0,
\label{eq:appinvalid}
\end{equation}
and
\begin{equation}
\mathrm{Cov}(I_t,I_{t+\ell}) = -p_w^2.
\label{eq:appinvalidcov}
\end{equation}

If $\ell \in S(w)$, then the joint event is feasible and
\begin{equation}
\begin{aligned}
\mathbb{P}(I_t=1,I_{t+\ell}=1)
&= \mathbb{P}(I_t=1)\mathbb{P}(I_{t+\ell}=1\mid I_t=1) \\
&= p_w q_w(\ell).
\end{aligned}
\label{eq:appvalid}
\end{equation}
Therefore
\begin{equation}
\mathrm{Cov}(I_t,I_{t+\ell})
=
p_w q_w(\ell)-p_w^2.
\label{eq:appvalidcov}
\end{equation}
This proves Equation~\eqref{eq:covlocal}. Summing over all pairs with the same shift gives
\begin{equation}
\begin{aligned}
\mathrm{Var}(N_w)
&= (T-m+1)p_w(1-p_w) \\
&\quad + 2\sum_{\ell=1}^{m-1}(T-m+1-\ell)\mathrm{Cov}(I_t,I_{t+\ell}) \\
&\quad + R_{\mathrm{long}},
\end{aligned}
\label{eq:applocalsum}
\end{equation}
where $R_{\mathrm{long}}$ collects shifts at least $m$. Dropping all nonpositive invalid overlap terms and all nonnegative contributions from $R_{\mathrm{long}}$ yields the lower bound in Equation~\eqref{eq:varcountlower}. This completes the proof.
\end{proof}

\section*{Proof of Theorem~\ref{thm:fisher}}\label{app:prooffisher}

Interpretation. This proof shows that a larger border score lowers retained information because suppression reacts to the larger burst scale of the response.

\begin{proof}
By definition,
\begin{equation}
I_w(\tau)
=
\mathbb{E}\!\left[\phi_{\tau}(a_w)^2 g_w^2\right].
\label{eq:appfisherdef}
\end{equation}
Since $g_w$ is independent of $a_w$ and $\mathbb{E}[g_w^2] < \infty$,
\begin{equation}
I_w(\tau)
=
\mathbb{E}[g_w^2]\,
\mathbb{E}\!\left[\phi_{\tau}(a_w)^2\right].
\label{eq:appfactor}
\end{equation}
Likewise,
\begin{equation}
I_w(\infty)
=
\mathbb{E}[g_w^2]\,
\mathbb{E}\!\left[\phi_{\infty}(a_w)^2\right].
\label{eq:appinfty}
\end{equation}
Because $\phi_{\infty}(u)=1$ by definition of no suppression, we obtain
\begin{equation}
I_w(\infty)=\mathbb{E}[g_w^2].
\label{eq:appinfty2}
\end{equation}
Hence
\begin{equation}
R_w(\tau)
=
\frac{I_w(\tau)}{I_w(\infty)}
=
\mathbb{E}\!\left[\phi_{\tau}(a_w)^2\right].
\label{eq:appratio}
\end{equation}
Using the scale model $a_w \overset{d}{=} b_w Z$ yields
\begin{equation}
R_w(\tau)
=
\mathbb{E}\!\left[\phi_{\tau}(b_w Z)^2\right].
\label{eq:appscaled}
\end{equation}
This is Equation~\eqref{eq:retentionexpect}.

To prove monotonicity, let $b_1 \leq b_2$. Since $\phi_{\tau}(u)$ is nonincreasing in $|u|$ and $|b_1 Z| \leq |b_2 Z|$ pointwise whenever $|Z|$ is fixed, we have
\begin{equation}
\phi_{\tau}(b_1 Z)^2 \geq \phi_{\tau}(b_2 Z)^2
\qquad \text{almost surely.}
\label{eq:appmono}
\end{equation}
Taking expectations gives
\begin{equation}
\mathbb{E}\!\left[\phi_{\tau}(b_1 Z)^2\right]
\geq
\mathbb{E}\!\left[\phi_{\tau}(b_2 Z)^2\right].
\label{eq:appmonoexp}
\end{equation}
Thus $R_w(\tau)$ is nonincreasing in $b_w$. Since $b_w=b_0(1+\beta_w)$ with $b_0>0$, $R_w(\tau)$ is also nonincreasing in $\beta_w$. If $\phi_{\tau}$ is strictly decreasing on a set of positive probability, then the inequality is strict for any $b_1 < b_2$ on a set of positive probability, which implies strict monotonicity. This completes the proof.
\end{proof}

\section*{Proof of Theorem~\ref{thm:gamma}}\label{app:proofgamma}

Interpretation. This proof shows that the compensation weight should grow exactly where retained information becomes small.

\begin{proof}
The objective in Equation~\eqref{eq:gammaopt} is separable across words. For a fixed word $w \in \mathcal{H}$ define
\begin{equation}
J_w(\gamma_w)
=
\omega_w \frac{c_w}{\gamma_w^2 R_w(\tau)}
+
\lambda_{\Gamma}\gamma_w^2,
\qquad \gamma_w>0.
\label{eq:appjw}
\end{equation}
Differentiating gives
\begin{equation}
J_w'(\gamma_w)
=
-2\omega_w \frac{c_w}{\gamma_w^3 R_w(\tau)}
+
2\lambda_{\Gamma}\gamma_w.
\label{eq:appjprime}
\end{equation}
Setting the derivative to zero yields
\begin{equation}
\lambda_{\Gamma}\gamma_w^4
=
\frac{\omega_w c_w}{R_w(\tau)}.
\label{eq:appstationary}
\end{equation}
Since $\gamma_w>0$, the unique stationary point is
\begin{equation}
\gamma_w^{\star}
=
\left(
\frac{\omega_w c_w}{\lambda_{\Gamma}R_w(\tau)}
\right)^{1/4}.
\label{eq:appclosedform}
\end{equation}
The second derivative is
\begin{equation}
J_w''(\gamma_w)
=
6\omega_w \frac{c_w}{\gamma_w^4 R_w(\tau)}
+
2\lambda_{\Gamma}
>
0,
\label{eq:appsecond}
\end{equation}
so the stationary point is the unique global minimizer. Finally, because $\gamma_w^{\star}$ is proportional to $R_w(\tau)^{-1/4}$, any decrease of $R_w(\tau)$ with $\beta_w$ implies a nondecreasing dependence of $\gamma_w^{\star}$ on $\beta_w$. If the decrease is strict, then the increase is strict. This completes the proof.
\end{proof}

\section*{Additional derivation for Markov continuation mass}\label{app:markov}

Interpretation. This derivation clarifies how the background Markov model enters the continuation mass and why the theory remains tied to genomic word structure.

For a valid shift $\ell \in S(w)$, the event $I_t(w)=1$ fixes the letters $w_1,\ldots,w_m$ on positions $t,\ldots,t+m-1$. The event $I_{t+\ell}(w)=1$ requires the suffix letters $w_{m-\ell+1},\ldots,w_m$ to be followed by the trailing block $w_{m-\ell+1},\ldots,w_m$ again. Under a first order Markov chain this gives
\begin{equation}
q_w(\ell)
=
\prod_{j=m-\ell+1}^{m-1}
P_{w_j,w_{j+1}},
\label{eq:appmarkovq}
\end{equation}
up to the shared overlap block that is already fixed by $I_t(w)=1$. Therefore the continuation mass can be estimated by
\begin{equation}
\widehat{\rho}_w
=
\sum_{\ell \in S(w)}
\prod_{j=m-\ell+1}^{m-1}
\widehat{P}_{w_j,w_{j+1}},
\label{eq:appmarkovrho}
\end{equation}
where $\widehat{P}$ is the empirical transition matrix estimated from the pretraining corpus. This appendix section is useful for implementation because it gives a direct route from tokenizer vocabulary to border scores without requiring labeled data. On heterogeneous genomes, $\widehat{\rho}_w$ should be interpreted cautiously when $w$ lies in repetitive or low-complexity regions because continuation mass can be inflated by tandem or interspersed repeats rather than by a single functional motif.

\section*{Implementation mapping from theory to trainer}\label{app:implementation}

Interpretation. This section specifies how tokenizer vocabulary, border scores, and compensation are instantiated in the released trainer, resolving the notational gap between vocabulary-aligned $\Gamma$ and the deployed LoRA branch.

\paragraph{Tokenizer-to-word alignment.}
Each tokenizer token $w \in \mathcal{V}$ is mapped to a DNA string (BPE merge or $k$-mer in DNABERT-2). For token $w$ of nucleotide length $m(w)$, the overlap set $S(w)$ and border polynomial $B_w(z)$ are computed on that string. The continuation mass $\widehat{\rho}_w$ follows Equation~\eqref{eq:appmarkovrho} with $\widehat{P}$ estimated from a background corpus. The border score is $\beta_w = \widehat{\rho}_w / (1-\widehat{\rho}_w)$ when $\widehat{\rho}_w < 1$.

\paragraph{Controlled metadata score.}
On synthetic border-aware splits, each sample stores a scalar \texttt{border\_score} in CSV metadata aligned with the constructed difficulty of that sample. During training, the score is batch-normalized and passed to Equation~\eqref{eq:implcomp} with $\lambda_{\mathrm{comp}}$, $c_{\min}$, and $c_{\max}$ from the YAML config.

\paragraph{External sequence-estimated score.}
When metadata is unavailable, a label-free score is computed from the input sequence (center-window JSD over $k$-mers), then train-quantile normalized to reduce saturation before compensation. The same Equation~\eqref{eq:implcomp} path is used; only the score source changes.

\paragraph{Low-rank branch.}
For each targeted linear layer $W$, LoRA factors $A$ and $B$ are attached as in standard LoRA. GERM-BO replaces the additive update $y = Wx + \mathrm{scale}\cdot B(Ax)$ with $y = Wx + \mathrm{scale}\cdot B(Ax \cdot c(x))$, where $c(x)$ is broadcast across valid token positions in a sequence. No separate $\Gamma$ matrix is stored; the theoretical diagonal $\gamma_w$ is realized through $c(x)$ derived from $\beta_w$ or metadata.

\paragraph{Hyperparameter correspondence.}
Config field \texttt{compensation\_strength} equals $\lambda_{\mathrm{comp}}$ in Equation~\eqref{eq:implcomp}. Fields \texttt{compensation\_clip\_min} and \texttt{compensation\_clip\_max} equal $c_{\min}$ and $c_{\max}$. LoRA \texttt{rank} and \texttt{alpha} are $r$ and the LoRA scaling numerator. The theoretical $\lambda_{\Gamma}$, $\tau$, and $\alpha$ in Equations~\eqref{eq:gammaopt} and~\eqref{eq:powerlaw} are not exposed as trainer knobs in the current release.

\end{document}